In recent years, microservice systems have gained significant popularity in the development of cloud-based applications, owing to their numerous advantages such as resource flexibility, a loosely coupled architecture, and lightweight deployment. However, failures are inevitable in microservice systems due to their inherent complexity. A failure in one service can propagate across the system, affecting many other services and resulting in the degradation of the system availability. This, in turn, leads to poor user experience and incurs huge economic losses. For instance, it has been reported that a one-hour downtime on Amazon.com could potentially cost up to 100 million USD~\cite{Chen2019incienttriage, Chen2020incidentmanagment}. Therefore, system operators must closely monitor the systems, checking key run-time information to promptly detect failures as soon as they occur, and then proceed to identify the failures' root causes and troubleshoot them. However, in practice, the complexity of microservice systems and the large volume of monitoring data make these tasks especially challenging.

Metric-based anomaly detection and root cause analysis (RCA) for microservice systems have been extensively studied in recent years \cite{Soldani2022rcasurvey, Azam2022rcd, Li2022Circa, Xin2023CausalRCA, chen2022adaptive, siffer2017anomaly, wu2020microrca, Jinjin2018Microscope, yu2021microrank, liu2021microhecl,lee2023eadro}. Given a set of metrics data, anomaly detection techniques aim to detect whether there exist anomalies and consequently, failures within the microservice system~\cite{chen2022adaptive, lee2023eadro}. If a failure is detected, the RCA module is then triggered to locate the root cause of the failure \cite{lee2023eadro, Azam2022rcd, Li2022Circa}. The RCA module aims to address two fundamental questions: (1) which services are the root causes, and (2) what specific issues are causing that failure (e.g. high CPU utilization, memory leak, or network congestion). Some RCA approaches use hypothesis testing or a statistical analysis method to analyze the time series metrics data to identify the candidate root causes for the detected anomalies \cite{Shan2019Ediagnosis, Li2022Circa}. Some RCA methods construct topology graphs using the information provided by the monitoring systems, such as the microservice status, the interaction between the services, and the interaction traces, to facilitate root cause analysis \cite{wu2020microrca, liu2021microhecl}. Multiple recent RCA methods \cite{Xin2023CausalRCA, Azam2022rcd, Jinjin2018Microscope, Meng2020Microcause} use different causal discovery methods \cite{Runge2019PCMCI, Jaber2020PsiFCI, Spirtes1995fci} to derive the causal relationships among the services and metrics from the multivariate time series metrics data and employ graph centrality algorithms like random walk \cite{Jinjin2018Microscope}, PageRank \cite{Xin2023CausalRCA,wu2021microdiag} or Depth-First Search (DFS) \cite{Chen2014Causeinfer} to infer the root causes.

Despite being closely related, existing RCA works typically either treat the anomaly detection tasks independently or rely on overly simplistic anomaly detection techniques~\cite{lee2023eadro, Soldani2022rcasurvey, Li2022Circa, Azam2022rcd}. For example, MicroRCA \cite{wu2020microrca} and MicroDiag \cite{wu2021microdiag} employ the simple BIRCH clustering technique~\cite{zhang1996birch} for detecting anomalies. MicroScope \cite{Jinjin2018Microscope} and Ms-Rank \cite{Ma2019Msrank} use a basic three-sigma rule of thumb as their anomaly detection method, known as N-Sigma. 
There are various commercial monitoring platforms, notable platforms including DataDog \cite{datadog} and Dynatrace \cite{dynatrace}, which rely on simplistic %threshold-based 
anomaly detection techniques \cite{dynatrace1, datadog2} for univariate time series data. 
Most RCA research works~\cite{Jinjin2018Microscope, Wang2018cloudranger, Ma2019Msrank, wu2020microrca, Ma2020Automap, wu2021microdiag, Li2022Circa, Azam2022rcd, Xin2023CausalRCA} focus solely on identifying the root cause of the failure whilst assuming the existence of an anomaly detection module that can accurately detect failures and trigger the RCA module when failures are detected. Some of these works, such as CIRCA \cite{Li2022Circa} and RCD \cite{Azam2022rcd}, specifically require certain information from the anomaly detection module, in particular, the failure occurrence time. However, they assume this information is already known accurately. Thus, it remains unclear whether, when combined with existing anomaly detection methods that may provide imprecise information, these approaches are still effective in localizing the root cause. 

% Dynatrace AD: If you want Dynatrace to suggest an appropriate threshold, select Use suggested threshold. Dynatrace will calculate a suggested threshold based on the last seven days of historical data. -> threshold-based anomaly detection
% https://www.dynatrace.com/news/blog/metric-events-set-up-anomaly-detection-based-on-your-business-needs/
% https://www.dynatrace.com/platform/artificial-intelligence/auto-baselining/
% https://www.dynatrace.com/news/blog/why-averages-suck-and-percentiles-are-great/
% Percentiles are also great for tuning and giving your optimizations a particular goal -> they use percentiles for optimization purpose

% Dynatrace RCA: After Davis identifies a problem.., IT USES ALL monitored transactions (distributed traces) to identify interdependencies....
% https://docs.dynatrace.com/docs/platform/davis-ai/problem-and-root-cause/root-cause-analysis#a-context-aware-approach-for-the-detec
% Whenever an event is triggered on a component, the AI root-cause analysis automatically collects all the transactions (distributed traces) along the horizontal stack. The analysis automatically continues when the horizontal stack shows that a called service is also marked unhealthy. With each hop on the horizontal stack, the vertical technology stack is also collected and analyzed for unhealthy states.


% https://docs.dynatrace.com/docs/platform/davis-ai/problem-and-root-cause/root-cause-analysis#expand--details-and-example


% Datadog AD: https://www.datadoghq.com/blog/introducing-anomaly-detection-datadog/

% Datadog RCA:  Watchdog RCA requires the use of APM (traces)
% https://docs.datadoghq.com/watchdog/rca/ 

% Datadog RCA: We are currently working to expand our coverage to other infrastructure issues, such as CPU saturation and memory leaks. https://www.datadoghq.com/blog/datadog-watchdog-automated-root-cause-analysis/


% \add{\hh{Think of where to put this one later} The novelty of our proposed BARO lies in its use of MultivariateBOCPD for anomaly detection and the introduction of a novel RobustScorer hypothesis testing technique for RCA. Our extensive investigations into commercial tools (e.g. DataDog and Dynatrace) revealed no mentions of MultivariateBOCPD or hypothesis testing, further validating our assertion that BARO represents a novel contribution that has yet to be introduced in both academic literature and commercial tools.}

% \lp{We are trying to address these two comments:
% - an explanation of why and how the approach is different from the commercial ones;
% - Authors should try to look at several commercial offerings from DataDog, Dynatrace, etc. and they will find similarities to the proposed framework.} \hh{Add 1-2 sentences saying there are some commercial tools and list their limitations}


In this work, we introduce BARO, an end-to-end approach for anomaly detection and root cause analysis for microservice systems based on metrics data, which are multivariate time series data. BARO includes a Multivariate Bayesian Online Change Point Detection module for detecting anomalies. It also includes a novel RobustScorer module, which is a nonparametric statistical hypothesis testing technique and less sensitive to the accuracy of the anomaly detection, for robustly identifying the root causes. BARO offers several advantages. Firstly, it follows an unsupervised learning approach, eliminating the need for labelled data and enabling direct application without the requirement of such labelled data. Secondly, it does not rely on operational knowledge (e.g., service call graphs) or causal graphs, making it suitable for large-scale evolving systems where acquiring such operational knowledge or causal graphs for numerous services is difficult \cite{Azam2022rcd, Li2022Circa, Wang2018cloudranger}. Finally, BARO is nonparametric, scale-equivalent, and rotation-invariant, making it applicable to a wide range of systems. We comprehensively evaluate BARO against various state-of-the-art approaches on three popular benchmark microservice systems. Our experimental results demonstrate that BARO consistently surpasses the state-of-the-art methods. Additionally, we analyze the sensitivity of the RCA methods against their parameters to show the robustness of our method.

In summary, our major contributions are as follows: 
\begin{itemize}
    \item We propose a new end-to-end approach for anomaly detection and root cause analysis in microservice systems based on multivariate time series metrics data. In particular, we propose to use the Multivariate Bayesian Online Change Point Detection technique to detect anomalies, and a novel nonparametric statistical hypothesis testing technique for %robustly 
    accurately identifying root causes of microservice systems' failures.
    \item We conduct extensive experiments on three popular benchmark microservice systems. Our experimental results demonstrate that BARO consistently outperforms state-of-the-art approaches in both anomaly detection and root cause analysis.
    \item We perform a comprehensive sensitivity analysis, evaluating the performance of all the RCA methods w.r.t. their parameters. Our experimental results show that BARO is significantly more robust against important parameters% compared to baseline methods.
    , e.g., the anomaly detection time, compared to baseline methods.
\end{itemize}

